The $\Sigma_1$ theory of non-structural subtyping consists of the true Boolean combinations
of existential and universal sentences over subtype inequalities. We prove its
decidability over arbitrary, regular, and finite trees. For a binary covariant constructor
with bottom and top, truth of an existential sentence is \NExp-complete and truth of a
universal sentence is \coNExp-complete. Boolean combinations can be decided in polynomial
time with parallel queries to an \NExp\ oracle.

The first step reduces the theory to one problem. A universal sentence is the negation of an
existential one, and an existential sentence holds exactly when, for some clause of its
disjunctive normal form, the positive inequalities do not entail the disjunction of the
negated ones. Everything thus rests on \emph{disjunctive entailment},
$\Gamma\entails x_1\le y_1\lor\cdots\lor x_k\le y_k$. Unlike plain entailment ($k=1$), the
disjuncts cannot be decided one at a time, since a countermodel must refute all of them in
one solution. We read the witness paths as words of a finite automaton that also tests pairs
of positions, within and across the words, and show that a tuple of singly exponential total
length suffices: we extend all words until some stop and the others continue periodically,
shorten the finite parts together on one trie of reversed blocks, and cut the periodic words
at lengths given by strict difference constraints. The results extend to a constructor of
any arity and variance, and decidability to the covariant signature with bottom alone.
